\subsection*{调试大量使用 MMA 的核函数}

MMA 路径出现错误结果时，先核对数据布局：操作数是行优先还是列优先，右侧矩阵在内存中是 \co@[N, K]@ 还是 \co@[K, N]@。再检查分块索引，确认 \co@.at@ 或 \co@chunkat@ 选择了正确的 \co@block_m@、\co@block_n@ 和 K 维切片。

常见问题包括实际数据布局与 \co@mma.row.row@ 的声明不一致，或分块尺寸与 MMA 所需形状不匹配。使用带 swizzle 的装入方式时，布局不一致还可能产生规律性的错误，例如每隔十六个元素只有一个正确。记录错误位置的规律，有助于缩小排查范围。
